\documentclass{article}
\usepackage[margin=1in]{geometry}

\usepackage{hyperref}
\usepackage{url}

\usepackage[utf8]{inputenc}
\usepackage[T1]{fontenc}
\usepackage{booktabs}
\usepackage{amsfonts}
\usepackage{nicefrac}
\usepackage{microtype}
\usepackage{xcolor}
\usepackage{fancyvrb}

\usepackage{amsmath}
\usepackage{amssymb}
\usepackage{mathtools}
\usepackage{amsthm}

\usepackage[inline]{enumitem}
\setitemize{noitemsep,topsep=0pt,parsep=0pt,partopsep=0pt,leftmargin=*}
\setenumerate{noitemsep,topsep=0pt,parsep=0pt,partopsep=0pt,leftmargin=*}

\usepackage[capitalize,noabbrev]{cleveref}
\crefname{section}{\S}{\S\S}
\crefname{figure}{Fig.}{Figs.}
\crefname{algorithm}{Alg.}{}
\crefname{equation}{Eq.}{Eq.}

\usepackage{latexsym}
\usepackage{cancel}
\usepackage{inconsolata}
\usepackage{amssymb}
\usepackage{mleftright}
\usepackage{xspace}

\usepackage{minted}
\setminted{linenos, firstnumber=last, escapeinside=@@, numbersep=4pt}

\usepackage{multicol}
\usepackage{setspace}
\usepackage{framed}

\usepackage{tikz}
\usetikzlibrary{tikzmark}
\usetikzlibrary{arrows}
\usetikzlibrary{arrows.meta}
\usetikzlibrary{shapes}
\usetikzlibrary{automata,positioning}
\tikzset{
  CenterFix/.style={baseline={(current bounding box.center)}},
}

\usepackage{graphicx}
\usepackage{bbm}
\usepackage{dsfont}
\usepackage{caption}
\usepackage{subcaption}
\usepackage{soul}
\usepackage{float}
\usepackage{stmaryrd}

\usepackage[normalem]{ulem}

\hypersetup{
  colorlinks=true,
  linkcolor=blue!35!black,
  citecolor=blue!35!black,
  urlcolor=blue!35!black
}

\theoremstyle{plain}
\newtheorem{theorem}{Theorem}[section]
\newtheorem{proposition}[theorem]{Proposition}
\newtheorem{lemma}[theorem]{Lemma}
\newtheorem{corollary}[theorem]{Corollary}
\newtheorem{definition}[theorem]{Definition}
\newtheorem{remark}[theorem]{Remark}

%%%%%%%%%%%%%%%%%%%%%%%%%%%%%%%%%%%%%%%%%%%%%%%%%%%%%%%%%%
% MACROS (from paper/macros.tex)
%%%%%%%%%%%%%%%%%%%%%%%%%%%%%%%%%%%%%%%%%%%%%%%%%%%%%%%%%%

\colorlet{MacroColor}{black}
\newcommand{\mymacro}[1]{#1}

\renewcommand{\Pr}[2][]{\mathop{\mathrm{Pr}}_{#1}\mleft[#2\mright]}

\newcommand{\prooftext}[1]{\text{\color{black!50}#1}}

\newcommand{\defn}[1]{\textbf{#1}}
\newcommand{\defeq}{\mathrel{\overset{\raisebox{-0.25ex}{\textnormal{\tiny def}}}{=}}}

\newcommand{\card}[1]{\left|#1\right|}
\newcommand{\set}[1]{\left\{ #1 \right\}}
\newcommand{\indicator}[1]{\mathbbm{1}\mleft\{#1\mright\}}

\newcommand{\bigo}{{\mymacro{ \mathcal{O}}}}
\newcommand{\bigO}[1]{{\mymacro{ \bigo\mleft(#1\mright)}}}

\newcommand{\kleene}[1]{#1^*}
\newcommand{\kleeneplus}[1]{#1^+}

\DeclareMathOperator*{\argmax}{argmax}
\DeclareMathOperator*{\argmin}{argmin}

\newcommand{\wsp}{\textvisiblespace}

%%%%%%%%%%%%%%%%%%%%%%%%%%%%%%%%%%%%%%%%%%%%%%%%%%%%%%%%%%
% PAPER-SPECIFIC NOTATION
%%%%%%%%%%%%%%%%%%%%%%%%%%%%%%%%%%%%%%%%%%%%%%%%%%%%%%%%%%

% Cover, quotient, remainder
\newcommand{\preCover}{{\color{TokenColor}\mathcal{P}}}
\newcommand{\Quotient}{{\color{TokenColor}\mathcal{Q}}}
\newcommand{\Remainder}{{\color{TokenColor}\mathcal{R}}}
\newcommand{\Quotientp}{\Quotient'}
\newcommand{\Remainderp}{\Remainder'}
\newcommand{\maxQuotient}{{\color{TokenColor}\mathcal{C}}}

\newcommand{\eos}{{\mymacro{\textsc{eos}}}}

\newcommand{\basisSet}[1]{\langle #1\rangle}

\definecolor{CharacterColor}{RGB}{98,115,19}
\definecolor{WordColor}{RGB}{102,0,204}
\definecolor{TokenColor}{RGB}{140,10,89}
\definecolor{TargetColor}{RGB}{98,115,19}
\definecolor{SourceColor}{RGB}{140,10,89}
\definecolor{StormBlue}{RGB}{42,90,120}

\newcommand{\tokenfont}[2][]{%
    \ensuremath{%
        \stackrel{\raisebox{-0.5ex}{{\color{TokenColor}\strut\texttt{#2}}}}%
        {\text{\scriptsize \raisebox{-0.5ex}{{\color{black!50}\texttt{\strut#1}}}}}%
    }%
}
\newcommand{\charfont}[1]{{\color{CharacterColor}\strut\texttt{#1}}}
\newcommand{\tfont}[1]{{\color{TokenColor}\texttt{#1}}}

\newcommand{\transFnSym}[0]{f}
\newcommand{\transFn}{\ensuremath{{\color{CharacterColor}\transFnSym}}\xspace}

% Source / target symbols and strings
\newcommand{\srcSym}{{\color{TokenColor}x}}
\newcommand{\srcSyma}{{\color{TokenColor}a}}
\newcommand{\srcSymp}{{\color{TokenColor}x^\prime}}
\newcommand{\srcStr}{\ensuremath{{\color{TokenColor}\boldsymbol{x}}}\xspace}
\newcommand{\srcStrp}[0]{{\color{TokenColor}\srcStr^\prime}}
\newcommand{\srcStrpp}[0]{{\color{TokenColor}\srcStr^{\prime\prime}}}

\newcommand{\tgtSym}{{\color{CharacterColor}y}}
\newcommand{\tgtSymb}{{\color{CharacterColor}b}}
\newcommand{\tgtSymp}{{\color{CharacterColor}y^\prime}}
\newcommand{\tgtSympp}{{\color{CharacterColor}y^{\prime\prime}}}
\newcommand{\tgtSymScore}{{\color{CharacterColor}\hat{y}}}

\newcommand{\tgtStr}{{\color{CharacterColor}\boldsymbol{y}}}
\newcommand{\tgtStrp}{{\color{CharacterColor}\boldsymbol{y}^\prime}}
\newcommand{\tgtStrpp}{{\color{CharacterColor}\boldsymbol{y}^{\prime\prime}}}

% Language models
\newcommand{\pSource}{\mymacro{p_{\srcAlphabet}}}
\newcommand{\pTarget}{\mymacro{p_{\tgtAlphabet}}}
\newcommand{\prefixLanguageOp}[1]{\overrightarrow{#1}}
\newcommand{\prefixSource}{{\mymacro{\prefixLanguageOp{\pSource}}}}
\newcommand{\prefixTarget}{{\mymacro{\prefixLanguageOp{\pTarget}}}}

% Alphabets
\newcommand{\srcAlphabet}[0]{{\color{TokenColor}\ensuremath{\mathcal{X}}}\xspace}
\newcommand{\tgtAlphabet}[0]{{\color{CharacterColor}\ensuremath{\mathcal{Y}}}\xspace}
\newcommand{\srcStrings}[0]{\ensuremath{\srcAlphabet^*}\xspace}
\newcommand{\tgtStrings}[0]{\ensuremath{\tgtAlphabet^*}\xspace}

% Transducer
\newcommand{\transST}{{\color{CharacterColor}\texttt{f}}}

% States
\newcommand{\States}{\mymacro{\mathsf{S}}}
\newcommand{\Transitions}{\mymacro{\mathsf{T}}}
\newcommand{\InitialStates}[0]{\mymacro{\mathsf{I}}}
\newcommand{\FinalStates}{\mymacro{\mathsf{F}}}

\newcommand{\state}{\mymacro{s}}
\newcommand{\statep}{\mymacro{s^\prime}}
\newcommand{\statepp}{\mymacro{s^{\prime\prime}}}

% Cover algorithm
\newcommand{\queue}[0]{\texttt{q}}
\newcommand{\candidates}[0]{\texttt{C}}
\newcommand{\frontier}{\mathcal{F}}
\newcommand{\frontierp}{\frontier^\prime}

\newcommand{\powerstate}{S}
\newcommand{\powerstatep}{\powerstate^{\prime}}

\newcommand{\threshold}{\tau}

\newcommand{\sem}[1]{\llbracket #1 \rrbracket}

\newcommand{\forceStart}[2]{#1_{[#2]}}

\newcommand{\buffer}{{\color{CharacterColor}\boldsymbol{b}}}
\newcommand{\bufferp}{{\color{CharacterColor}\boldsymbol{b}'}}

% Minted function names
\newcommand{\mintedFunc}[1]{\mbox{\texttt{{\color{orange!50!black}#1}}}}
\newcommand{\myMintedFunc}[1]{\mbox{\texttt{{\color{StormBlue}#1}}}}

\newcommand{\continuous}[0]{\myMintedFunc{is\_cylinder}}
\newcommand{\discontinuous}[0]{\myMintedFunc{is\_member}}
\newcommand{\candidate}[0]{\myMintedFunc{is\_live}}
\newcommand{\build}[0]{\myMintedFunc{build}}
\newcommand{\prune}[0]{\myMintedFunc{prune}}
\newcommand{\decompose}[0]{\myMintedFunc{decompose}}

\newcommand{\runFrontier}[0]{\myMintedFunc{run}}
\newcommand{\nextFrontier}[0]{\myMintedFunc{step}}
\newcommand{\closureFrontier}[0]{\myMintedFunc{closure}}

\newcommand{\reachableOutputs}[0]{\myMintedFunc{reachable\_outputs}}
\newcommand{\reachableInputs}[0]{\myMintedFunc{reachable\_inputs}}
\newcommand{\isExactMember}[0]{\myMintedFunc{is\_exact\_member}}
\newcommand{\decomposeNext}[0]{\myMintedFunc{decompose\_next}}

\newcommand{\ipUniversal}[0]{\myMintedFunc{fast\_univ\_filter}}
\newcommand{\ipUniversalStates}[0]{U}

\newcommand{\projUniv}[0]{\myMintedFunc{proj\_pos}}

\newcommand{\ProjA}{\mymacro{\mathsf{proj}}_{\srcAlphabet}}

\newcommand{\pdict}{\mymacro{\overline{p}}}
\newcommand{\pdictNew}{\mymacro{\overline{p}_{\texttt{new}}}}

\newcommand{\Visited}{\mymacro{V}}

\newcommand{\Cache}{\mymacro{\texttt{cache}}}

%%%%%%%%%%%%%%%%%%%%%%%%%%%%%%%%%%%%%%%%%%%%%%%%%%%%%%%%%%

\title{Source-String Incremental Decomposition\\with Multi-Symbol Expansion}
\author{}
\date{}

\begin{document}
\maketitle

\section{Overview}

We present an incremental decomposition algorithm that operates directly on
\emph{source strings} $\srcStr \in \srcStrings$, maintaining a frontier of
$(\state, \buffer)$ pairs internally but exposing the decomposition as sets of
source strings.
The algorithm has two key components:
\begin{enumerate}
\item $\decompose(\tgtStr)$: a memoized, incremental BFS that computes
  $(\Remainder(\tgtStr), \Quotient(\tgtStr))$ by refining the parent
  decomposition $(\Remainder(\tgtStr_{<N}), \Quotient(\tgtStr_{<N}))$.
\item $\decomposeNext(\tgtStr)$: a single-worklist BFS that computes
  $\Remainder(\tgtStr\tgtSym)$ and $\Quotient(\tgtStr\tgtSym)$ for \emph{all}
  $\tgtSym \in \tgtAlphabet$ simultaneously, then populates the $\decompose$
  cache to avoid redundant BFS for subsequent calls.
\end{enumerate}
The next-symbol distribution $\prefixTarget(\cdot \mid \tgtStr)$ follows
directly from $\decomposeNext$, giving a much simpler algorithm than the
general \texttt{fast\_next\_dist}.

\paragraph{Notation.}
We follow the notation of the main paper.
$\transST$ is a functional FST with source alphabet $\srcAlphabet$ and target
alphabet $\tgtAlphabet$.
$\tgtStr \in \tgtStrings$ is a target string (prefix) with $N \defeq |\tgtStr|$.
$\Remainder(\tgtStr)$ and $\Quotient(\tgtStr)$ denote the remainder and
quotient of the prefix decomposition
$\transFn^{-1}(\tgtStr\tgtStrings) = \Remainder(\tgtStr) \cup
\Quotient(\tgtStr) \cdot \srcStrings$.
$\pSource$ is a language model over source strings, and
$\prefixSource(\srcStr) \defeq \sum_{\srcStrp \in \srcStrings}
\pSource(\srcStr\srcStrp)$ is its prefix probability.
A \defn{frontier} $\frontier \defeq \{(\state, \buffer)\}$ is a set of
(FST state, output buffer) pairs.


\section{Frontier Computation}
\label{sec:frontier}

The frontier $\runFrontier_\tgtStr(\srcStr)$ computes the set of $(\state,
\buffer)$ pairs reachable by reading source string $\srcStr$ from the initial
states, filtering to buffers compatible with $\tgtStr$ (\cref{fig:frontier}).
Compatibility means $\tgtStr$ and $\buffer$ agree on their shared prefix:
$\tgtStr_{<k} = \buffer_{<k}$ where $k = \min(N, |\buffer|)$.
This allows buffers shorter than $\tgtStr$ (output not yet produced) or longer
(output already past target --- still live for prefix matching).

The frontier supports two recursion strategies (\cref{fig:frontier}):
\begin{itemize}
\item \textbf{Source recursion:}
  $\runFrontier_\tgtStr(\srcStr) =
  \nextFrontier_\tgtStr(\runFrontier_\tgtStr(\srcStr_{<N}), \srcSym_N)$.
  Extends one source symbol with target fixed.
\item \textbf{Target recursion:}
  $\runFrontier_{\tgtStr\tgtSym}(\srcStr) =
  \{(\state, \buffer) \in \runFrontier_\tgtStr(\srcStr) :
  |\buffer| \le N \lor \buffer_{N+1} = \tgtSym\}$.
  Derives from parent target by filtering
  (\cref{prop:frontier-containment}).
\end{itemize}
Both are memoized on $(\srcStr, \tgtStr)$.  Source recursion is the default.
Target recursion is used opportunistically when
$\runFrontier_\tgtStr(\srcStr)$ is already cached --- typical in
$\decomposeNext$, where we move from target $\tgtStr$ to target
$\tgtStr\tgtSym$.  This avoids recomputing the entire source-string chain for
the new target.

\begin{proposition}[Frontier containment]
\label{prop:frontier-containment}
For any source string $\srcStr$ and target symbol $\tgtSym$,
\[
\runFrontier_{\tgtStr\tgtSym}(\srcStr) = \{(\state, \buffer) \in
\runFrontier_\tgtStr(\srcStr) : |\buffer| \le N \text{ or } \buffer_{N+1} =
\tgtSym\}.
\]
\end{proposition}
\begin{proof}
Buffer filtering in $\nextFrontier$ and $\closureFrontier$ keeps only
$(\state, \buffer)$ pairs where $\buffer$ and the target agree on their shared
prefix.
Since FST arcs only append to the buffer (never backtrack), once
$\buffer_{N+1} = \tgtSymp \ne \tgtSym$ at some step, all descendant buffers
also have position $N{+}1 = \tgtSymp$ and would be filtered out by the
$\tgtStr\tgtSym$ filter.
Therefore, filtering $\runFrontier_\tgtStr(\srcStr)$ at the end is equivalent
to filtering at each step.

The $(\supseteq)$ direction holds because $\tgtStr\tgtSym$-compatible elements
are also $\tgtStr$-compatible (the former is strictly more restrictive).
The $(\subseteq)$ direction holds because elements that are
$\tgtStr$-compatible but \emph{not} $\tgtStr\tgtSym$-compatible (i.e.,
$|\buffer| > N$ and $\buffer_{N+1} \ne \tgtSym$) cannot produce descendants
that become $\tgtStr\tgtSym$-compatible, since buffers only grow.
\end{proof}

\begin{figure}[H]
\centering
\begin{minipage}[t]{.35\linewidth}
\begin{minted}{python}
def @$\runFrontier_\tgtStr(\srcSym_1 \cdots \srcSym_N)$@:  # memoize
  if @$N = 0$@:
    return @$\closureFrontier_\tgtStr(\InitialStates \times \{ \varepsilon \})$@
  # Target recursion (if parent cached)
  if @$|\tgtStr| > 0$@ and @$\runFrontier_{\tgtStr_{<N}}(\srcStr)$@ in cache:
    @$\tgtSym \gets \tgtStr_N$@; @$M \gets N{-}1$@
    return @$\{(\state, \buffer) \in \runFrontier_{\tgtStr_{<N}}(\srcStr)$@
           @$\colon |\buffer| \le M \lor \buffer_{M+1} = \tgtSym \}$@
  # Source recursion (default)
  return @$\nextFrontier_\tgtStr(\runFrontier_\tgtStr(\srcStr_{<N}), \srcSym_N)$@
\end{minted}
\end{minipage}
\hfill
\begin{minipage}[t]{.30\linewidth}
\begin{minted}{python}
def @$\nextFrontier_\tgtStr(\frontier, \srcSym)$@:
  @$\frontierp \gets \emptyset$@
  for @$(\state, \buffer)$@ in @$\frontier$@:
    for @$(\_ \xrightarrow{\srcSym:\tgtSym} \statep)\in \Transitions(\state, \srcSym)$@:
      @$\bufferp  \gets \buffer\tgtSym$@
      if @$\bufferp \preceq \tgtStr \lor \tgtStr \preceq \bufferp$@:
        @$\frontierp$@.add(@$(\statep, \bufferp)$@)
  return @$\closureFrontier_\tgtStr(\frontierp)$@
\end{minted}
\end{minipage}
\hfill
\begin{minipage}[t]{.30\linewidth}
\begin{minted}{python}
def @$\closureFrontier_\tgtStr(\frontier)$@:
  @$\frontierp \gets \frontier$@
  @$\queue \gets \textsc{Queue}(\frontier)$@
  while @$|\queue| > 0$@:
    @$(\state, \buffer) \gets$@ @$\queue$@.pop()
    for @$(\_ \xrightarrow{\varepsilon:\tgtSymp} \statep)\in \Transitions(\state, \varepsilon)$@:
      @$\bufferp  \gets \buffer\tgtSymp$@
      if @$\bufferp \preceq \tgtStr \lor \tgtStr \preceq \bufferp$@:
        if @$(\statep, \bufferp) \notin \frontierp$@:
          @$\frontierp$@.add(@$(\statep, \bufferp)$@)
          @$\queue$@.add(@$(\statep, \bufferp)$@)
  return @$\frontierp$@
\end{minted}
\end{minipage}
\caption{Frontier computation.
$\runFrontier_\tgtStr(\srcStr)$ supports two recursion strategies:
source recursion via $\nextFrontier$ (extending $\srcStr$) and
target recursion via filtering (extending $\tgtStr$,
\cref{prop:frontier-containment}).
Target recursion is used when the parent-target frontier is already cached,
avoiding a full chain recomputation.
Buffer filtering discards $(\state, \buffer)$ pairs where $\buffer$ is
incompatible with $\tgtStr$; sound because FSTs only append to the buffer.
We write $\buffer\tgtSym$ as shorthand for $\buffer$ when $\tgtSym = \varepsilon$,
and $\buffer \cdot (\tgtSym)$ otherwise.}
\label{fig:frontier}
\end{figure}


\section{Frontier-Based Checks}
\label{sec:checks}

The decomposition relies on three checks, each defined in terms of the frontier
(\cref{fig:checks}).  These are identical to the main paper's
\texttt{LazyCachedDecomp}.
We do not use the $\mintedFunc{combined\_universal}$
shortcut from the paper; see the comparison section below for details.

\begin{figure}[H]
\centering
\begin{minipage}[t]{.58\linewidth}
\begin{minted}{python}
class Incremental:
  @$({\_, \srcAlphabet, \_, \Transitions, \InitialStates, \FinalStates}) \gets \transST$@
\end{minted}
\begin{minted}{python}
  def @$\continuous(\srcStr, \tgtStr)$@:  # memoize
    @$N \gets |\tgtStr|$@
    def @$\projUniv(\frontierp)$@:
      return @$\{(\state, \bufferp_{<N}) \mid (\state, \bufferp) \in \frontierp,$@
             @$\bufferp \preceq \tgtStr \lor \tgtStr \preceq \bufferp \}$@
    @$\powerstate \gets \projUniv(\runFrontier_\tgtStr(\srcStr))$@
    @$\Visited \gets \{ \powerstate \}$@
    @$\queue \gets \textsc{Queue}(\{\powerstate\})$@
    while @$\queue$@:
      @$\powerstate \gets \queue$@.pop()
      # Accepting: some pair has matched
      # target and is accepting
      if @$\neg \exists\, (\state, \buffer) \in \powerstate \colon \buffer \succeq \tgtStr \land \state \in \FinalStates$@:
        return False
      # Complete: every input leads to a
      # non-empty successor
      for @$\srcSymp \in \srcAlphabet$@:
        @$\powerstatep \gets \projUniv(\nextFrontier_\tgtStr(\powerstate, \srcSymp))$@
        if @$\powerstatep = \emptyset$@: return False
        if @$\powerstatep \notin \Visited$@:
          @$\Visited$@.add(@$\powerstatep$@); @$\queue$@.push(@$\powerstatep$@)
    return True
\end{minted}
\end{minipage}
\hfill
\begin{minipage}[t]{.38\linewidth}
\begin{minted}{python}
  def @$\discontinuous(\srcStr, \tgtStr)$@:
    return @$\exists (\state, \buffer) \in$@
      @$\runFrontier_\tgtStr(\srcStr) \colon$@
      @$\buffer \succeq \tgtStr \land \state \in \FinalStates$@
\end{minted}
\begin{minted}{python}
  def @$\candidate(\srcStr, \tgtStr)$@:
    return @$\runFrontier_\tgtStr(\srcStr) \ne \emptyset$@
\end{minted}
\begin{minted}{python}
  def @$\isExactMember(\srcStr, \tgtStr)$@:
    return @$\exists (\state, \buffer) \in$@
      @$\runFrontier_\tgtStr(\srcStr) \colon$@
      @$\buffer = \tgtStr \land \state \in \FinalStates$@
\end{minted}
\begin{minted}{python}
  def @$\prune(\candidates)$@:
    return @$\{\srcStr \in \candidates$@
      @$\mid \prefixSource(\srcStr) > 0 \}$@
\end{minted}
\begin{minted}{python}
  def @$\reachableInputs(\srcStr, \tgtStr)$@:  # memoize
    @$N \gets |\tgtStr|$@
    return @$\{\srcSymp \mid (\state, \buffer) \in \runFrontier_\tgtStr(\srcStr),$@
      @$(\state \xrightarrow{\srcSymp:\tgtSym} \statep) \in \Transitions(\state, \srcSymp),$@
      @$\buffer\tgtSym \preceq \tgtStr \lor \tgtStr \preceq \buffer\tgtSym \}$@
\end{minted}
\begin{minted}{python}
  def @$\reachableOutputs(\srcStr, \tgtStr)$@:
    @$N \gets |\tgtStr|$@
    @$\frontier \gets \runFrontier_\tgtStr(\srcStr)$@
    return @$\{\bufferp_{N+1} \mid$@
      @$(\state, \buffer) \in \frontier,$@
      @$(\state \xrightarrow{\_:\tgtSym} \statep) \in \Transitions(\state),$@
      @$\bufferp = \buffer\tgtSym, |\bufferp| > N,$@
      @$\bufferp_{<N} = \tgtStr_{<N} \}$@
    # over-approx: also yields @$\tgtSym$@
    # where buffer has not yet reached
    # position @$N$@; filtered by @$\candidate$@ later
\end{minted}
\end{minipage}
\caption{Frontier-based checks on $\texttt{Incremental}$.
\textbf{Left:} $\continuous(\srcStr, \tgtStr)$ performs powerset BFS to verify
that every continuation of $\srcStr$ produces output starting with $\tgtStr$.
Buffers are truncated to length $N$ via $\projUniv$ to keep the powerset
state space finite.
\textbf{Right:}
$\discontinuous$ checks if $\srcStr$ produces output covering $\tgtStr$ at a
final state.
$\candidate$ checks if $\srcStr$ has any live frontier state.
$\isExactMember$ checks if $\transFn(\srcStr) = \tgtStr$ exactly.
$\prune$ drops source strings with zero LM prefix probability.
$\reachableInputs$ returns source symbols leading to target-compatible
successors (used in $\decompose$ to prune the extension loop).
$\reachableOutputs$ returns target symbols $\tgtSym$ reachable at position
$N{+}1$ of the output (used in $\decomposeNext$ to prune per-symbol checks).}
\label{fig:checks}
\end{figure}


\section{Incremental Decomposition}
\label{sec:decompose}

The $\decompose$ method (\cref{fig:decompose}) is a memoized, incremental BFS
that computes $(\Remainder(\tgtStr), \Quotient(\tgtStr))$.
For the empty target $\varepsilon$, the BFS is seeded with $\{\varepsilon\}$.
For $\tgtStr = \tgtStr_{<N}\tgtSym_N$, the BFS is seeded with
$\Quotient(\tgtStr_{<N}) \cup \Remainder(\tgtStr_{<N})$ from the parent
decomposition.
Parent quotient strings may lose universality under the new target symbol;
parent remainder strings may gain new extensions.
Each source string is classified:
\begin{itemize}
\item \textbf{Cylinder} ($\continuous$): $\srcStr$ and all its extensions
  produce output starting with $\tgtStr$.  Added to $\Quotient$;
  no further extension needed (the entire subtree is absorbed).
\item \textbf{Member} ($\discontinuous$): $\srcStr$ itself produces output
  starting with $\tgtStr$.  Added to $\Remainder$; extensions are still
  explored (some may be cylinders or members of their own).
\item \textbf{Extension}: for each reachable source symbol $\srcSymp$,
  $\srcStr\srcSymp$ is added to the candidate set if it has a live frontier
  ($\candidate$).
\end{itemize}
The BFS terminates because $\continuous$ absorbs entire subtrees ($\Quotient$
grows monotonically in depth), $\candidate$ filters dead prefixes, and $\prune$
filters LM-dead prefixes.

\begin{figure}[H]
\centering
\begin{minted}{python}
class Incremental:  # continued
  def @$\decompose(\tgtStr)$@:  # memoize
    @$N \gets |\tgtStr|$@
    if @$N = 0$@:
      @$\queue \gets \{\varepsilon\}$@
    else:
      @$(\Remainderp, \Quotientp) \gets \decompose(\tgtStr_{<N})$@
      @$\queue \gets \Quotientp\cup\Remainderp$@
    @$(\Remainder, \Quotient) \gets (\emptyset, \emptyset)$@
    while @$|\queue| > 0$@:
      @$\candidates \gets \emptyset$@
      for @$\srcStr$@ in @$\queue$@:
        if @$\continuous(\srcStr, \tgtStr)$@:
          @$\Quotient$@.add(@$\srcStr$@); continue        # absorb entire subtree
        if @$\discontinuous(\srcStr, \tgtStr)$@:
          @$\Remainder$@.add(@$\srcStr$@)
        # not a cylinder @$\Rightarrow$@ must explore extensions
        for @$\srcSymp \in \reachableInputs(\srcStr, \tgtStr)$@:
          @$\candidates$@.add(@$\srcStr\srcSymp$@)
      @$\queue \gets \prune(\candidates)$@
    return @$(\Remainder, \Quotient)$@
\end{minted}
\caption{Incremental decomposition.
$\decompose(\tgtStr)$ is memoized: extending to $\tgtStr\tgtSym$ starts from
$\Quotient(\tgtStr) \cup \Remainder(\tgtStr)$ rather than $\varepsilon$.
The extension loop uses $\reachableInputs$ rather than iterating over all of
$\srcAlphabet$, pruning source symbols whose arcs produce target-incompatible
buffers.}
\label{fig:decompose}
\end{figure}


\section{Multi-Symbol Expansion}
\label{sec:decompose-next}

The $\decomposeNext$ method (\cref{fig:decompose-next}) computes
$\Remainder(\tgtStr\tgtSym)$ and $\Quotient(\tgtStr\tgtSym)$ for all
$\tgtSym \in \tgtAlphabet$ simultaneously, as well as the exact preimage
$\{\srcStr : \transFn(\srcStr) = \tgtStr\}$.
This is the key efficiency improvement over calling $\decompose(\tgtStr\tgtSym)$
independently for each $\tgtSym$.

The BFS is seeded with $\Quotient(\tgtStr) \cup \Remainder(\tgtStr)$.
For each source string $\srcStr$ in the worklist:
\begin{enumerate}
\item \textbf{Exact preimage:} if $\isExactMember(\srcStr, \tgtStr)$,
  add $\srcStr$ to the preimage set.
\item \textbf{Per-symbol classification:} for each
  $\tgtSym \in \reachableOutputs(\srcStr, \tgtStr)$, check if $\srcStr$ is a
  cylinder or member for $\tgtStr\tgtSym$.
  By functionality of $\transST$, at most one $\tgtSym$ can be a cylinder
  (see below).
\item \textbf{Absorption:} if a cylinder was found for some $\tgtSym$, the
  entire subtree rooted at $\srcStr$ is absorbed --- skip extension for
  \emph{all} symbols.
\item \textbf{Extension:} otherwise, extend by all reachable source symbols.
\end{enumerate}

Three optimizations:
\begin{itemize}
\item \textbf{Target recursion in $\runFrontier$ (\cref{sec:frontier}):}
  The call $\continuous(\srcStr, \tgtStr\tgtSym)$ internally triggers
  $\runFrontier_{\tgtStr\tgtSym}(\srcStr)$.
  Because $\decomposeNext$ has already computed
  $\runFrontier_\tgtStr(\srcStr)$ (via $\decompose(\tgtStr)$), the frontier
  for the extended target is derived by filtering rather than recomputing the
  entire source-string chain (\cref{prop:frontier-containment}).
  This is transparent: $\continuous$ is called unchanged; the optimization
  lives entirely inside $\runFrontier$.
\item \textbf{Non-cylinder shortcut:}
  Strings in $\Remainder(\tgtStr)$ were not cylinders for $\tgtStr$,
  hence cannot be cylinders for the strictly more restrictive $\tgtStr\tgtSym$.
  The expensive powerset universality check is skipped for these seeds.
\item \textbf{Cache population:}
  After the BFS, the computed $\Remainder(\tgtStr\tgtSym)$ and
  $\Quotient(\tgtStr\tgtSym)$ are written into $\decompose$'s memoization
  cache.
  A subsequent $\decompose(\tgtStr\tgtSym)$ call (e.g., from
  $\decomposeNext(\tgtStr\tgtSym)$) hits the cache instead of re-running its
  own BFS.
\end{itemize}

\begin{proposition}[Cylinder uniqueness]
\label{prop:cylinder-uniqueness}
For a functional FST $\transST$, a source string $\srcStr$ can be a cylinder
for at most one target symbol $\tgtSym$.
\end{proposition}
\begin{proof}
Suppose $\srcStr$ is a cylinder for both $\tgtSym$ and $\tgtSymp$ with $\tgtSym
\ne \tgtSymp$.
Being a cylinder for $\tgtSym$ means every extension of $\srcStr$ produces
output starting with $\tgtStr\tgtSym$; likewise for $\tgtSymp$.
For a functional FST, each input has a unique output, so
$\transFn^{-1}(\tgtStr\tgtSym\tgtStrings) \cap
\transFn^{-1}(\tgtStr\tgtSymp\tgtStrings) = \emptyset$.
But $\srcStr$ and all its extensions lie in both preimages --- a contradiction.
\end{proof}

\begin{figure}[H]
\centering
\begin{minted}{python}
  def @$\decomposeNext(\tgtStr)$@:
    @$(\Remainder, \Quotient) \gets \decompose(\tgtStr)$@
    @$\texttt{quotients} \gets \{\}$@; @$\texttt{remainders} \gets \{\}$@        # @$\tgtSym \mapsto$@ set of @$\srcStr$@
    @$\texttt{preimage} \gets \emptyset$@
    # Remainder strings were not cylinders for @$\tgtStr$@, hence not for any @$\tgtStr\tgtSym$@.
    @$\texttt{non\_cyl} \gets \Remainder$@
    @$\queue \gets \Remainder \cup \Quotient$@
    while @$|\queue| > 0$@:
      @$\candidates \gets \emptyset$@
      for @$\srcStr$@ in @$\queue$@:
        if @$\isExactMember(\srcStr, \tgtStr)$@:
          @$\texttt{preimage}$@.add(@$\srcStr$@)               # @$\transFn(\srcStr) = \tgtStr$@ exactly
        # Check all reachable @$\tgtSym$@'s for this @$\srcStr$@ in one pass.
        # At most one @$\tgtSym$@ can be a cylinder (\cref{prop:cylinder-uniqueness}).
        @$\texttt{cyl\_sym} \gets \texttt{None}$@
        for @$\tgtSym$@ in @$\reachableOutputs(\srcStr, \tgtStr)$@:
          if @$\srcStr \notin \texttt{non\_cyl}$@ and @$\continuous(\srcStr, \tgtStr\tgtSym)$@:
            assert @$\texttt{cyl\_sym} = \texttt{None}$@      # functionality
            @$\texttt{quotients}[\tgtSym]$@.add(@$\srcStr$@); @$\texttt{cyl\_sym} \gets \tgtSym$@; continue
          if @$\discontinuous(\srcStr, \tgtStr\tgtSym)$@:
            @$\texttt{remainders}[\tgtSym]$@.add(@$\srcStr$@)
        if @$\texttt{cyl\_sym} \ne \texttt{None}$@: continue   # absorbed @$\Rightarrow$@ skip extension for all @$\tgtSym$@
        for @$\srcSymp \in \reachableInputs(\srcStr, \tgtStr)$@:
          @$\candidates$@.add(@$\srcStr\srcSymp$@)
      @$\queue \gets \prune(\candidates)$@
    # Populate @$\decompose$@'s cache for all children
    for @$\tgtSym \in \texttt{quotients}.\texttt{keys}() \cup \texttt{remainders}.\texttt{keys}()$@:
      @$\Cache[\decompose, \tgtStr\tgtSym] \gets (\texttt{remainders}[\tgtSym], \texttt{quotients}[\tgtSym])$@
    return @$(\texttt{remainders}, \texttt{quotients}, \texttt{preimage})$@
\end{minted}
\caption{Multi-symbol expansion.
$\decomposeNext(\tgtStr)$ computes $\Remainder(\tgtStr\tgtSym)$ and
$\Quotient(\tgtStr\tgtSym)$ for all $\tgtSym \in \tgtAlphabet$ in a single
BFS pass.
The $\continuous(\srcStr, \tgtStr\tgtSym)$ call triggers
$\runFrontier_{\tgtStr\tgtSym}(\srcStr)$, which is computed cheaply via
target recursion (\cref{prop:frontier-containment}) since the parent frontier
$\runFrontier_\tgtStr(\srcStr)$ is already cached.
The non-cylinder shortcut avoids the universality check entirely for
seeds in $\Remainder(\tgtStr)$.
After the BFS, the cache is populated so that subsequent
$\decompose(\tgtStr\tgtSym)$ calls are free.
The exact preimage is collected as a byproduct: if $\srcStr$ is a cylinder for
some $\tgtStr\tgtSym$, every extension produces output strictly longer than
$\tgtStr$, so no extension can be an exact preimage of $\tgtStr$.}
\label{fig:decompose-next}
\end{figure}


\section{Scoring}
\label{sec:scoring}

Given the decomposition, the scoring functions (\cref{fig:scoring}) are
straightforward.
The prefix probability $\prefixTarget(\tgtStr)$ sums LM prefix probabilities
over $\Quotient$ and LM string probabilities over $\Remainder$.
The next-symbol distribution $\prefixTarget(\cdot \mid \tgtStr)$ is computed
via $\decomposeNext$, which provides the per-symbol decompositions and exact
preimage in a single pass.

\begin{figure}[H]
\centering
\begin{minipage}[t]{.48\linewidth}
\begin{minted}{python}
  def prefix_prob(@$\tgtStr$@):
    @$(\Remainder, \Quotient) \gets \decompose(\tgtStr)$@
    return @$\sum_{\srcStr \in \Quotient} \prefixSource(\srcStr) + \sum_{\srcStr \in \Remainder} \pSource(\srcStr)$@
\end{minted}
\begin{minted}{python}
  def prob(@$\tgtStr$@):
    @$(\_, \_, \texttt{preimage}) \gets \decomposeNext(\tgtStr)$@
    return @$\sum_{\srcStr \in \texttt{preimage}} \pSource(\srcStr)$@
\end{minted}
\end{minipage}
\hfill
\begin{minipage}[t]{.48\linewidth}
\begin{minted}{python}
  def next_dist(@$\tgtStr$@):
    @$(\texttt{remainders}, \texttt{quotients}, \texttt{preimage}) \gets \decomposeNext(\tgtStr)$@
    @$\pdict \gets \{\}$@
    for @$\tgtSym \in \texttt{quotients}.\texttt{keys}() \cup \texttt{remainders}.\texttt{keys}()$@:
      for @$\srcStr \in \texttt{quotients}[\tgtSym]$@:
        @$\pdict[\tgtSym]$@ += @$\prefixSource(\srcStr)$@
      for @$\srcStr \in \texttt{remainders}[\tgtSym]$@:
        @$\pdict[\tgtSym]$@ += @$\pSource(\srcStr)$@
    for @$\srcStr \in \texttt{preimage}$@:
      @$\pdict[\eos]$@ += @$\pSource(\srcStr)$@
    return @$\mintedFunc{normalize}(\pdict)$@
\end{minted}
\end{minipage}
\caption[Scoring functions]{%
\textbf{Left:} $\mintedFunc{prefix\_prob}(\tgtStr) = \prefixTarget(\tgtStr)$ and
$\mintedFunc{prob}(\tgtStr) = \pTarget(\tgtStr)$.
\textbf{Right:} $\mintedFunc{next\_dist}(\tgtStr) = \prefixTarget(\cdot \mid \tgtStr)$
via $\decomposeNext$.
Compared to $\mintedFunc{fast\_next\_dist}$, this is simpler: there is no
per-element expansion loop and no $\bot$ / catching-up logic.
Target recursion in $\runFrontier$ (\cref{prop:frontier-containment}) avoids
the frontier recomputation.
The entire complexity is absorbed by $\decomposeNext$, which classifies all
symbols in one pass.
All computations are in the log domain; we write sums for clarity.
$\prefixSource$ and $\pSource$ are memoized so that
$\prefixSource(\cdot \mid \cdot)$ takes $\mathcal{O}(1)$ time.}
\label{fig:scoring}
\end{figure}

For reference, the na\"ive next-symbol distribution computes
$\mintedFunc{prefix\_prob}(\tgtStr\tgtSym)$ independently for each $\tgtSym$:

\begin{figure}[H]
\centering
\begin{minted}{python}
  def next_dist_bruteforce(@$\tgtStr$@):
    @$\pdict \gets \{\}$@
    @$Z \gets \mintedFunc{prefix\_prob}(\tgtStr)$@
    for @$\tgtSymp \in \tgtAlphabet$@:
      @$\pdict[\tgtSymp] \gets \mintedFunc{prefix\_prob}(\tgtStr\tgtSymp) / Z$@
    @$\pdict[\eos] \gets \mintedFunc{prob}(\tgtStr) / Z$@
    return @$\pdict$@
\end{minted}
\caption{Reference implementation.
Correct but redundant: each $\mintedFunc{prefix\_prob}(\tgtStr\tgtSym)$ call
triggers its own $\decompose$ BFS.
Used as a test oracle for $\mintedFunc{next\_dist}$.}
\label{fig:bruteforce}
\end{figure}


\section{Correctness}

\begin{proposition}
$\mintedFunc{next\_dist}(\tgtStr) = \mintedFunc{next\_dist\_bruteforce}(\tgtStr)$.
\end{proposition}

\begin{proof}
It suffices to show that $\decomposeNext(\tgtStr)$ produces the same
$\Remainder(\tgtStr\tgtSym)$ and $\Quotient(\tgtStr\tgtSym)$ as calling
$\decompose(\tgtStr\tgtSym)$ independently.

Both $\decompose(\tgtStr\tgtSym)$ and $\decomposeNext(\tgtStr)$ are seeded
with $\Quotient(\tgtStr) \cup \Remainder(\tgtStr)$.
The former classifies each $\srcStr$ using $\continuous(\srcStr,
\tgtStr\tgtSym)$ and $\discontinuous(\srcStr, \tgtStr\tgtSym)$, extending
non-cylinders by $\reachableInputs(\srcStr, \tgtStr\tgtSym)$.
The latter performs the same checks for all $\tgtSym$ simultaneously.

There are three potential differences to verify:
\begin{enumerate}
\item \textbf{Reachable symbols:} $\decomposeNext$ only checks
  $\tgtSym \in \reachableOutputs(\srcStr, \tgtStr)$ rather than all of
  $\tgtAlphabet$.
  By soundness of $\reachableOutputs$, any $\tgtSym$ with
  $\Quotient(\tgtStr\tgtSym) \ne \emptyset$ or
  $\Remainder(\tgtStr\tgtSym) \ne \emptyset$ is included.

\item \textbf{Extension symbols:} $\decomposeNext$ extends by
  $\reachableInputs(\srcStr, \tgtStr)$ while $\decompose(\tgtStr\tgtSym)$
  extends by $\reachableInputs(\srcStr, \tgtStr\tgtSym)$.
  The former is a superset of the latter (compatibility with $\tgtStr$ is weaker
  than compatibility with $\tgtStr\tgtSym$), so no reachable extensions are
  missed.

\item \textbf{Cylinder absorption:} when $\srcStr$ is a cylinder for $\tgtSym$,
  $\decomposeNext$ skips extension for \emph{all} symbols.
  This is correct by
  \cref{prop:cylinder-uniqueness}: $\srcStr$ cannot be a cylinder for any other
  $\tgtSymp \ne \tgtSym$.
  Furthermore, since $\srcStr \in \Quotient(\tgtStr\tgtSym)$, every extension
  of $\srcStr$ produces output starting with $\tgtStr\tgtSym$ --- none can
  contribute to $\Quotient(\tgtStr\tgtSymp)$ or $\Remainder(\tgtStr\tgtSymp)$
  for $\tgtSymp \ne \tgtSym$, and none can contribute additional members to
  $\Remainder(\tgtStr\tgtSym)$ beyond what the cylinder already absorbs.
\end{enumerate}

\noindent
The non-cylinder shortcut is also correct: if $\srcStr \in \Remainder(\tgtStr)$,
then $\srcStr$ was not a cylinder for $\tgtStr$.
Since $\tgtStr\tgtSym$ is strictly more restrictive than $\tgtStr$ (being a
cylinder for $\tgtStr\tgtSym$ implies being a cylinder for $\tgtStr$),
$\srcStr$ cannot be a cylinder for any $\tgtStr\tgtSym$.
\end{proof}



\section{Relationship to Existing Algorithms}

\paragraph{Comparison with \texttt{AbstractIncrDecomp}.}
The $\decompose$ method is identical to
\texttt{AbstractIncrDecomp}.$\decompose$ from the main paper, with one
refinement: the extension loop iterates over
$\reachableInputs(\srcStr, \tgtStr) \subseteq \srcAlphabet$ rather than all of
$\srcAlphabet$, pre-filtering source symbols whose arcs produce
target-incompatible buffers.

\paragraph{Comparison with \texttt{fast\_next\_dist}.}
The $\mintedFunc{next\_dist}$ function achieves the same result as
\texttt{fast\_next\_dist} but with a fundamentally different structure.
\texttt{fast\_next\_dist} interleaves symbol classification with source-string
expansion in a single loop, requiring $\bot$ markers, a
$\mintedFunc{combined\_universal}$ helper, and catching-up logic for beams
whose buffers lag behind the target.
In contrast, $\mintedFunc{next\_dist}$ separates the concerns cleanly:
$\decomposeNext$ handles all the BFS and classification, and
$\mintedFunc{next\_dist}$ simply reads off the results.
The cost of this simplicity is that $\decomposeNext$ always materializes the
full decomposition $(\Remainder(\tgtStr\tgtSym), \Quotient(\tgtStr\tgtSym))$
for every reachable $\tgtSym$, whereas \texttt{fast\_next\_dist} can sometimes
resolve a symbol early (via $\mintedFunc{combined\_universal}$) without full
expansion.

\paragraph{Comparison with \texttt{combined\_universal}.}
The paper's $\mintedFunc{combined\_universal}$ is a \emph{sufficient condition}
for cylinder status: it partitions the frontier into ``scored'' states
(buffer past target, committed to $\tgtSym$) and ``unscored'' states
(buffer = target), checks universality of their union on plain FST states
(no buffer tracking), and verifies that unscored arcs produce only $\tgtSym$
or $\varepsilon$.
When it succeeds, it avoids \emph{both} the frontier recomputation and the
full powerset BFS of $\continuous$.
When it fails, the algorithm falls through to full expansion.

Our algorithm addresses the frontier recomputation cost via
\textbf{target recursion} (\cref{sec:frontier}): by
\cref{prop:frontier-containment}, the $\tgtStr\tgtSym$ frontier is derived by
filtering the cached $\tgtStr$ frontier ($\mathcal{O}(|\frontier|)$),
avoiding the chain recomputation.
This is exact, transparent, and always applies.
We do not currently use a universality shortcut analogous to
$\mintedFunc{combined\_universal}$; the full powerset BFS in $\continuous$
handles all cases.

\paragraph{Comparison with \texttt{LazyCachedDecomp}.}
The frontier-based checks ($\continuous$, $\discontinuous$, $\candidate$) and
frontier computation ($\runFrontier$, $\nextFrontier$, $\closureFrontier$) are
identical to \texttt{LazyCachedDecomp} from the main paper.
The only structural difference is that this algorithm operates on source strings
(passed to $\continuous(\srcStr, \tgtStr)$ etc.) while
\texttt{LazyCachedDecomp} sometimes operates on frontiers directly.
Internally, both compute the same frontier via $\runFrontier_\tgtStr(\srcStr)$.


\end{document}
